\documentclass[10pt,a4paper]{article}

% Loesungsfassung. Die Schuelerfassung
% arbeitsblatt-1-atome-und-ladungen.tex ist zeichengleich und
% unterscheidet sich ausschliesslich in der Zeile \solutionfalse / \solutiontrue.

\usepackage[T1]{fontenc}
\usepackage[utf8]{inputenc}
\usepackage[ngerman]{babel}
\usepackage[a4paper,margin=14mm]{geometry}
\usepackage{lmodern}
\usepackage[table]{xcolor}
\usepackage{tabularx}
\usepackage{array}
\usepackage{amsmath}
\usepackage{microtype}
\usepackage[most]{tcolorbox}
\usepackage{tikz}
\usetikzlibrary{arrows.meta,calc,positioning}

\newif\ifsolution
\solutiontrue

\pagestyle{empty}
\setlength{\parindent}{0pt}
\setlength{\parskip}{0pt}
\renewcommand{\familydefault}{\sfdefault}
\renewcommand{\arraystretch}{1.08}

\definecolor{Navy}{HTML}{17324D}
\definecolor{Blue}{HTML}{3B6EA5}
\definecolor{Sky}{HTML}{EAF3FA}
\definecolor{Mint}{HTML}{E8F5EF}
\definecolor{Gold}{HTML}{F3B33D}
\definecolor{Ink}{HTML}{1E2935}
\definecolor{Muted}{HTML}{617080}
\definecolor{Line}{HTML}{D8E1E8}

\color{Ink}

\newcommand{\sectiontitle}[1]{%
  \vspace{1.0mm}{\color{Blue}\bfseries\large #1}\par
  \vspace{0.3mm}{\color{Blue}\rule{\linewidth}{0.7pt}}\par
  \vspace{0.6mm}}

\newcommand{\parttitle}[2]{%
  \vspace{1.4mm}%
  \colorbox{Navy}{\parbox{\dimexpr\linewidth-2\fboxsep\relax}{%
    \vspace{0.9mm}\color{white}\bfseries #1\hfill{\small\mdseries #2}\vspace{0.9mm}}}%
  \par\vspace{0.8mm}}

% ---------------------------------------------------------------------------
% Arbeitsbereiche: in Schueler- und Loesungsfassung exakt gleich hoch
% ---------------------------------------------------------------------------
\newlength{\schreibzeile}
\setlength{\schreibzeile}{8.2mm}
\newcounter{zeilenzaehler}

\makeatletter
% #1 Anzahl der Schreiblinien, #2 Breite. Die leere hbox legt die Breite der
% umgebenden vbox fest; ohne sie waeren die \hrule-Linien null Punkt breit.
\newcommand{\zeilenstapel}[2]{%
  \hbox to #2{}%
  \setcounter{zeilenzaehler}{0}%
  \@whilenum\value{zeilenzaehler}<#1\do{%
    \stepcounter{zeilenzaehler}%
    \vskip\dimexpr\schreibzeile-0.6pt\relax
    {\color{Muted!55}\hrule height 0.6pt}}}
\makeatother

\newcommand{\schreibraum}[3][\linewidth]{%
  \parbox[t][\dimexpr#2\schreibzeile\relax][t]{#1}{%
    \ifsolution
      \vspace*{1.0mm}\color{Blue}#3%
    \else
      \zeilenstapel{#2}{#1}%
    \fi}}

\newcommand{\gap}[2][28mm]{%
  \underline{\makebox[#1][c]{\strut\ifsolution\textcolor{Blue}{\bfseries #2}\fi}}}

\newcommand{\kasten}{%
  \tikz[baseline=-0.75mm]{\draw[Muted,line width=.55pt] (0,0) rectangle (3.6mm,3.6mm);}}
\newcommand{\kastenhaken}{%
  \tikz[baseline=-0.75mm]{\draw[Muted,line width=.55pt] (0,0) rectangle (3.6mm,3.6mm);
    \draw[Blue,line width=.9pt] (0.7mm,1.9mm)--(1.5mm,0.8mm)--(3.0mm,3.1mm);}}
\newcommand{\mck}[1]{\ifsolution\ifnum#1=1\relax\kastenhaken\else\kasten\fi\else\kasten\fi}

% Begriffskaertchen fuer den Wortspeicher unter einem Lueckentext
\newcommand{\wortchip}[1]{\fcolorbox{Line}{white}{\strut\small #1}}

\newcommand{\ladung}[1]{%
  \tikz[baseline=-0.95mm]{\node[circle,draw=Blue,fill=Sky,line width=.6pt,inner sep=0pt,
    minimum size=4.4mm,font=\scriptsize\bfseries,text=Blue]{#1};}}

% Papierschnipsel: unregelmaessige Umrisse, damit sie nicht wie Rechtecke wirken
\newcommand{\schnipsel}[3]{%
  \begin{scope}[shift={(#1,#2)},rotate=#3]
    \draw[fill=white,draw=Muted,line width=.4pt]
      (0,0)--(2.4,0.6)--(4.3,-0.2)--(3.9,-2.0)--(1.7,-2.5)--(0.2,-1.7)--cycle;
  \end{scope}}

% Zwei geriebene Luftballons an Faeden. #1 Laengeneinheit, #2 Bildunterschrift.
% Wird in A1 und noch einmal verkleinert in A3 verwendet.
\newcommand{\ballonpaar}[2]{%
  \begin{tikzpicture}[x=#1,y=#1,>=Latex,font=\scriptsize]
    \path (-28,-23) rectangle (28,13);
    \draw[Muted,line width=.8pt] (-16,11)--(16,11);
    \draw[Muted,line width=.5pt] (-4,11)--(-15,3.4);
    \draw[Muted,line width=.5pt] (4,11)--(15,3.4);
    \draw[fill=Sky,draw=Blue,line width=.7pt] (-15,-6.4) ellipse (6.5 and 8.4);
    \draw[fill=Sky,draw=Blue,line width=.7pt] (-16.3,1.0)--(-13.7,1.0)--(-15,3.4)--cycle;
    \draw[fill=Sky,draw=Blue,line width=.7pt] (15,-6.4) ellipse (6.5 and 8.4);
    \draw[fill=Sky,draw=Blue,line width=.7pt] (13.7,1.0)--(16.3,1.0)--(15,3.4)--cycle;
    \draw[-{Latex[length=1.8mm]},Blue,line width=.7pt] (-21.8,-6.4)--(-27,-6.4);
    \draw[-{Latex[length=1.8mm]},Blue,line width=.7pt] (21.8,-6.4)--(27,-6.4);
    \node[anchor=north,align=center,font=\scriptsize] at (0,-16.5) {#2};
  \end{tikzpicture}}

\newcommand{\blattfuss}{%
  \vfill
  {\color{Line}\rule{\linewidth}{0.5pt}}\par\vspace{1.0mm}
  \begin{tabularx}{\linewidth}{@{}Xr@{}}
    {\small\color{Muted}Klasse 8 · Physik · Elektrischer Strom} &
    {\small\color{Muted}Arbeitsblatt 1 · Atome und Ladungen\ifsolution{} · Lösungen\fi{} · Seite \thepage{} von 3}
  \end{tabularx}}

\begin{document}

% ===========================================================================
% Seite 1
% ===========================================================================
\colorbox{Navy}{%
  \parbox{\dimexpr\linewidth-2\fboxsep\relax}{%
    \vspace{2.0mm}\color{white}
    \begin{tabularx}{\linewidth}{@{}Xr@{}}
      {\small\bfseries PHYSIK · KLASSE 8} &
      \colorbox{Gold}{\color{Navy}\small\bfseries\strut \ifsolution LÖSUNGEN\else ARBEITSBLATT\fi}
    \end{tabularx}
    \vspace{1.6mm}
    {\fontsize{18}{21}\selectfont\bfseries Atome und Ladungen}\par
    \vspace{0.5mm}{\large Arbeitsblatt 1 -- Elektrischer Strom}\vspace{2.0mm}
  }%
}

\vspace{1.6mm}
\colorbox{Sky}{\parbox{\dimexpr\linewidth-2\fboxsep\relax}{%
  \vspace{0.9mm}\textcolor{Blue}{\bfseries Ziel:} Du erklärst den Aufbau eines Atoms und beschreibst mit Ladungen, warum sich Körper anziehen oder abstoßen und was beim Stromfluss passiert.\vspace{0.9mm}}}

% ---------------------------------------------------------------------------
\sectiontitle{A1 · Beobachten}
Die Luftballons wurden vorher an den Haaren gerieben.

\vspace{1.6mm}
\begin{minipage}[t]{0.48\linewidth}
\centering
\begin{tikzpicture}[x=1mm,y=1mm,>=Latex,font=\scriptsize]
  \path (-10,-23) rectangle (36,13);
  \draw[fill=Sky,draw=Blue,line width=.7pt] (0,0) ellipse (8 and 10);
  \draw[fill=Sky,draw=Blue,line width=.7pt] (-1.5,-9.6)--(1.5,-9.6)--(0,-12.6)--cycle;
  % einige Schnipsel haften am Ballon, andere liegen daneben
  \schnipsel{7.4}{3.6}{-18}
  \schnipsel{7.8}{-1.0}{12}
  \schnipsel{6.4}{-5.2}{-6}
  \schnipsel{17.0}{5.0}{24}
  \schnipsel{20.0}{-2.0}{-12}
  \schnipsel{16.5}{-8.2}{8}
  \node[anchor=north,font=\scriptsize] (lschnipsel) at (25,-13.5) {Papierschnipsel};
  \draw[dashed,-{Latex[length=1.8mm]},Muted,line width=.5pt] (lschnipsel) -- (20.6,-4.4);
\end{tikzpicture}
\end{minipage}\hfill
\begin{minipage}[t]{0.48\linewidth}
\centering
\ballonpaar{1mm}{zwei geriebene Luftballons\\an Fäden}
\end{minipage}

\vspace{2.6mm}
\textbf{Beschreibe}, was du auf dem Bild beobachtest. \textbf{Erkläre} die Beobachtung mit Hilfe von Ladungen.

\vspace{1.2mm}
\schreibraum{4}{Bild 1: Der geriebene Luftballon zieht die Papierschnipsel an. Bild 2: Die beiden geriebenen Luftballons stoßen sich ab.\par
Durch das Reiben werden die Luftballons elektrisch geladen. Beide Luftballons sind gleich geladen, deshalb stoßen sie sich ab. Die Papierschnipsel sind nicht geladen und werden vom geladenen Luftballon angezogen.}

\vspace{3.0mm}

% ---------------------------------------------------------------------------
\sectiontitle{A2 · Atome aufbauen}
Fülle den Lückentext aus und beschrifte die Skizze.

\vspace{1.6mm}
Jedes Atom eines Stoffes besteht aus positiv geladenen \gap[30mm]{Protonen}, sowie elektrisch neutralen \gap[32mm]{Neutronen}, die zusammen den \gap[27mm]{Atomkern} bilden, und entsprechend vielen \gap[24mm]{negativ} geladenen \gap[30mm]{Elektronen}, die an \gap[24mm]{den Kern} gebunden sind.

\vspace{3.0mm}
\begin{minipage}[c]{0.60\linewidth}
\centering
\begin{tikzpicture}[x=1mm,y=1mm,font=\small]
  % gestrichelt: Bereich, in dem sich die Elektronen aufhalten
  \draw[Line,dashed,line width=.7pt] (0,0) circle (15);
  % Atomkern: Protonen und Neutronen ueberlappen sich zu einem Koerper
  \node[circle,fill=Muted,inner sep=0pt,minimum size=4.6mm] (neutron) at (2.1,2.1) {};
  \node[circle,fill=Muted,inner sep=0pt,minimum size=4.6mm] at (-2.1,-2.1) {};
  \node[circle,fill=Blue,text=white,inner sep=0pt,minimum size=4.6mm,font=\scriptsize\bfseries] (proton) at (-2.1,2.1) {$+$};
  \node[circle,fill=Blue,text=white,inner sep=0pt,minimum size=4.6mm,font=\scriptsize\bfseries] at (2.1,-2.1) {$+$};
  % Elektronen: deutlich kleiner, in unterschiedlichem Abstand innerhalb der Huelle
  \node[circle,fill=white,draw=Blue,line width=.5pt,text=Blue,inner sep=0pt,minimum size=2.6mm,font=\tiny\bfseries] (elektron) at (10,-8) {$-$};
  \node[circle,fill=white,draw=Blue,line width=.5pt,text=Blue,inner sep=0pt,minimum size=2.6mm,font=\tiny\bfseries] at (-5,8) {$-$};
  \node[anchor=east] (lproton) at (-24,15) {\gap[24mm]{Proton}};
  \node[anchor=west] (lneutron) at (24,15) {\gap[26mm]{Neutron}};
  \node[anchor=east] (lkern) at (-24,-15) {\gap[24mm]{Atomkern}};
  \node[anchor=west] (lelektron) at (24,-15) {\gap[26mm]{Elektron}};
  \draw[Muted,line width=.5pt] (lproton) -- (proton);
  \draw[Muted,line width=.5pt] (lneutron) -- (neutron);
  \draw[Muted,line width=.5pt] (lkern) -- (0,-3.0);
  \draw[Muted,line width=.5pt] (lelektron) -- (elektron);
\end{tikzpicture}
\end{minipage}\hfill
\begin{minipage}[c]{0.36\linewidth}
\colorbox{Mint}{\parbox{\dimexpr\linewidth-2\fboxsep\relax}{%
  \vspace{1.0mm}\small\textcolor{Blue}{\bfseries Größenordnung}\par
  Ein Mensch besteht aus etwa $7\cdot10^{27}$ Atomen.\vspace{1.0mm}}}

\vspace{2.0mm}
{\small Die Skizze zeigt ein Heliumatom mit 2 Protonen, 2 Neutronen und 2 Elektronen.\par
\vspace{1.0mm}
Die gestrichelte Linie ist keine Bahn. Sie zeigt den Bereich, in dem sich die Elektronen aufhalten.}
\end{minipage}

\blattfuss
\newpage

% ===========================================================================
% Seite 2
% ===========================================================================
\sectiontitle{A3 · Ladungen vergleichen}
Trage ein, was zwischen den beiden Ladungen passiert.

\vspace{1.4mm}
\begin{tabularx}{\linewidth}{@{}p{72mm}X@{}}
  \rowcolor{Sky}\textbf{Ladungen} & \textbf{Was passiert zwischen ihnen?} \\
  \ladung{$+$}~\ladung{$+$}\quad zwei positive Ladungen & \gap[52mm]{stoßen sich ab} \\
  \ladung{$-$}~\ladung{$-$}\quad zwei negative Ladungen & \gap[52mm]{stoßen sich ab} \\
  \ladung{$+$}~\ladung{$-$}\quad eine positive, eine negative Ladung & \gap[52mm]{ziehen sich an}
\end{tabularx}

\vspace{1.8mm}
\colorbox{Mint}{\parbox{\dimexpr\linewidth-2\fboxsep\relax}{%
  \vspace{1.0mm}\textcolor{Blue}{\bfseries Merke:} Gleichnamige Ladungen \gap[34mm]{stoßen sich ab}, ungleichnamige Ladungen \gap[34mm]{ziehen sich an}.\vspace{1.0mm}}}

\vspace{1.8mm}
\textbf{Erkläre} mit dieser Erkenntnis, wie die beiden Luftballons beim zweiten Bild aus Aufgabe A1 geladen sind.

\vspace{1.2mm}
\begin{minipage}[c]{0.60\linewidth}
\schreibraum{3}{Die beiden Luftballons stoßen sich ab. Sie tragen also gleichnamige Ladungen.\par
Beide sind beim Reiben gleich geladen worden: entweder beide positiv oder beide negativ.}
\end{minipage}\hfill
\begin{minipage}[c]{0.34\linewidth}
\centering
\ballonpaar{0.9mm}{Bild 2 aus Aufgabe A1}
\end{minipage}

\vspace{3.0mm}

% ---------------------------------------------------------------------------
\sectiontitle{A4 · Stromfluss und Stromrichtung}
Fülle den Merksatz mit den Begriffen aus dem Wortspeicher aus. Zeichne danach beide Richtungen in den Stromkreis ein.

\vspace{1.4mm}
\colorbox{Sky}{\parbox{\dimexpr\linewidth-2\fboxsep\relax}{%
  \vspace{1.0mm}\textcolor{Blue}{\bfseries Merke:} Stromfluss ist die Bewegung von \gap[30mm]{Ladungen}. In guten Leitern wie Kupfer sind das \gap[30mm]{Elektronen}.\par
  \vspace{1.6mm}
  Die Stromrichtung ist durch die Bewegungsrichtung einer positiven Ladung festgelegt. Sie zeigt vom \gap[26mm]{Pluspol} zum \gap[26mm]{Minuspol} und ist der Bewegungsrichtung der \gap[30mm]{Elektronen} entgegengesetzt.\vspace{1.0mm}}}

\vspace{1.6mm}
{\small\color{Muted}Wortspeicher:}~\wortchip{Elektronen (2x)}~\wortchip{Ladungen}~\wortchip{Minuspol}~\wortchip{Pluspol}

\vspace{2.0mm}
\begin{minipage}[c]{0.60\linewidth}
\centering
\begin{tikzpicture}[x=1mm,y=1mm,>=Latex,font=\scriptsize]
  \path (-22,-10) rectangle (84,42);
  % Leitung; die senkrechten Enden beruehren die beiden Batterieplatten
  \draw[Navy,line width=.9pt] (0,0)--(0,12) (0,18)--(0,30)--(62,30)--(62,20);
  \draw[Navy,line width=.9pt] (62,10)--(62,0)--(44,0);
  \draw[Navy,line width=.9pt] (26,0)--(0,0);
  % Batterie
  \draw[Blue,line width=1.6pt] (-6,18)--(6,18);
  \draw[Blue,line width=1.6pt] (-3.2,12)--(3.2,12);
  \node[Blue,font=\small\bfseries] at (10,19.2) {$+$};
  \node[Blue,font=\small\bfseries] at (10,11.4) {$-$};
  \node[anchor=east,font=\scriptsize] at (-8,15) {Batterie};
  % Lampe
  \draw[fill=Sky,draw=Navy,line width=.9pt] (62,15) circle (5);
  \draw[Navy,line width=.7pt] (58.5,11.5)--(65.5,18.5) (58.5,18.5)--(65.5,11.5);
  \node[anchor=west,font=\scriptsize] at (68,15) {Lampe};
  % geschlossener Schalter
  \draw[fill=Navy,draw=Navy] (26,0) circle (.9) (44,0) circle (.9);
  \draw[Navy,line width=.9pt] (26,0)--(35,2.6)--(44,0);
  \node[anchor=north,font=\scriptsize] at (35,-3.4) {Schalter};
  \ifsolution
    \draw[-{Latex[length=2.2mm]},Blue,line width=1.0pt] (16,32.6)--(46,32.6);
    \node[anchor=south,font=\scriptsize\bfseries,text=Blue] at (31,33.8) {Stromrichtung};
    \draw[{Latex[length=2.2mm]}-,Blue,dashed,line width=1.0pt] (16,27.4)--(46,27.4);
    \node[anchor=north,font=\scriptsize\bfseries,text=Blue] at (31,26.2) {Bewegung der Elektronen};
  \fi
\end{tikzpicture}
\end{minipage}\hfill
\begin{minipage}[c]{0.36\linewidth}
\begin{tcolorbox}[colback=Mint,colframe=Line,arc=1.2mm,boxrule=.6pt,left=2mm,right=2mm,top=1.2mm,bottom=1.2mm]
{\small\textcolor{Blue}{\bfseries Hilfen zu A4}\par
\vspace{0.9mm}
\textbf{1.} Welche Teilchen können sich in einem Draht überhaupt bewegen?\par
\vspace{0.9mm}
\textbf{2.} Welche Ladung haben diese Teilchen? Von welchem Pol werden sie angezogen, von welchem werden sie abgestoßen?}
\end{tcolorbox}
\end{minipage}

\blattfuss
\newpage

% ===========================================================================
% Seite 3
% ===========================================================================
\parttitle{Teil B · Prüfungsaufgaben}{ohne Hilfen bearbeiten}

\sectiontitle{B1 · Bausteine des Atoms \hfill 3 BE}
Nenne die drei Bausteine eines Atoms und gib jeweils ihre Ladung an.

\vspace{1.2mm}
\schreibraum{3}{Protonen sind positiv geladen. Neutronen sind elektrisch neutral, sie tragen keine Ladung. Elektronen sind negativ geladen.\par
Protonen und Neutronen bilden zusammen den Atomkern. Die Elektronen sind an den Kern gebunden.}

\sectiontitle{B2 · Ein geriebener Glasstab \hfill 3 BE}
Ein Glasstab gibt beim Reiben mit einem Tuch Elektronen ab. Gib an, wie er danach geladen ist, und begründe deine Antwort.

\vspace{1.2mm}
\schreibraum{3}{Der Glasstab hat danach weniger Elektronen als vorher, also weniger negative Ladungen.\par
Die positiven Ladungen der Protonen überwiegen jetzt. Der Glasstab ist deshalb positiv geladen.}

\sectiontitle{B3 · Stromrichtung im Kupferdraht \hfill 3 BE}
In einem Kupferdraht bewegen sich die Elektronen vom Minuspol zum Pluspol. Gib die Stromrichtung an und begründe.

\vspace{1.2mm}
\schreibraum{3}{Die Stromrichtung zeigt vom Pluspol zum Minuspol.\par
Sie ist durch die Bewegungsrichtung einer positiven Ladung festgelegt. Deshalb ist sie der Bewegungsrichtung der Elektronen entgegengesetzt.}

\vspace{2.4mm}
\hfill{\small\color{Muted}erreichte Punkte: \gap[20mm]{} von 9 BE}

\vspace{3.4mm}

% ---------------------------------------------------------------------------
\sectiontitle{B4 · Check zum Abschluss \hfill Ankreuzen}
Kreuze alle richtigen Antworten an. Es können mehrere Antworten richtig sein.

\vspace{1.6mm}
\begin{minipage}[t]{0.485\linewidth}
\textbf{1. Welche Aussagen über Atome sind richtig?}\par
\vspace{1.0mm}
\mck{1}~Protonen sind positiv geladen.\par
\vspace{0.6mm}\mck{0}~Neutronen sind negativ geladen.\par
\vspace{0.6mm}\mck{1}~Elektronen sind an den Kern gebunden.\par
\vspace{0.6mm}\mck{1}~Der Kern besteht aus Protonen und Neutronen.

\vspace{3.0mm}
\textbf{3. Wie ist die Stromrichtung festgelegt?}\par
\vspace{1.0mm}
\mck{0}~Sie zeigt vom Minuspol zum Pluspol.\par
\vspace{0.6mm}\mck{1}~Sie zeigt vom Pluspol zum Minuspol.\par
\vspace{0.6mm}\mck{0}~Sie zeigt in Bewegungsrichtung der Elektronen.
\end{minipage}\hfill
\begin{minipage}[t]{0.485\linewidth}
\textbf{2. Was bewegt sich beim Stromfluss im Kupferdraht?}\par
\vspace{1.0mm}
\mck{1}~Ladungen\par
\vspace{0.6mm}\mck{1}~Elektronen\par
\vspace{0.6mm}\mck{0}~Protonen\par
\vspace{0.6mm}\mck{0}~der Atomkern

\vspace{3.0mm}
\textbf{4. Zwei Körper ziehen sich an. Was gilt?}\par
\vspace{1.0mm}
\mck{1}~Sie sind ungleichnamig geladen.\par
\vspace{0.6mm}\mck{0}~Sie sind gleichnamig geladen.\par
\vspace{0.6mm}\mck{1}~Einer ist positiv, einer negativ geladen.
\end{minipage}

\blattfuss

\end{document}
